\subsection{Regularization}

PROPOSITION 19. --- Let G be a locally compact group, \(\beta\) a relatively invariant positive measure \(\neq 0\) on G, B a base for the filter of neighborhoods of \(e\) in G, consisting of compact neighborhoods. For every V \(\in\) B, let \(f_{V}\) be a continuous function \(\geqslant 0\) on G, with support contained in V, such that \(\int f_{V} d\beta = 1\) . If \(\mu\) is a measure on G then, in \(\mathcal{M}(G)\) equipped with the topology of compact convergence in \(\mathcal{K}(G)\) ,

\[
\mu = \lim _ {\mathrm{V}} (\mu * f _ {\mathrm{V}}) \cdot \beta = \lim _ {\mathrm{V}} (f _ {\mathrm{V}} * \mu) \cdot \beta ,
\]

the limit being taken with respect to the section filter of B.

For the topology of compact convergence in \(\mathcal{C}(\mathrm{G})\) , \(f_{V} \cdot \beta\) tends to \(\varepsilon_{e}\) with respect to the section filter of B ( \(\S2\) , No.~7, Cor. 1 of Lemma 4). Therefore \(\mu = \lim_{\mathrm{V}} \mu * (f_{\mathrm{V}} \cdot \beta) = \lim_{\mathrm{V}} (f_{\mathrm{V}} \cdot \beta) * \mu\) in \(\mathcal{M}(\mathrm{G})\) equipped with the topology of compact convergence in \(\mathcal{K}(\mathrm{G})\) ( \(\S3\) , No.~3, Cor. of Prop. 12).

Remarks. --- 1) We thus see that every measure on G is the limit of measures admitting a continuous density with respect to every Haar measure (for the topology indicated in Prop. 19 and a fortiori for the vague topology).

\begin{enumerate}
\def\labelenumi{\arabic{enumi})}
\setcounter{enumi}{1}
\tightlist
\item
  If G is metrizable, B can be taken to be a sequence \((\mathrm{V}_{n})\) of neighborhoods. Then \(\mu\) is the limit of the sequence of measures \((\mu * f_{\mathrm{V}_{n}}) \cdot \beta\) with continuous densities. *If G is a real Lie group, the \(f_{V_{n}}\) can be taken to be infinitely differentiable; we shall see later on that the densities \(\mu * f_{V_{n}}\) are then infinitely differentiable.*
\end{enumerate}

PROPOSITION 20. --- We conserve the hypotheses and notations of Prop. 19. Let \(p \in [1, +\infty[\) and \(g \in \mathrm{L}^p(\mathbf{G}, \beta)\) . Then

\[
g = \lim _ {\mathrm{V}} g * ^ {\beta} f _ {\mathrm{V}} = \lim _ {\mathrm{V}} f _ {\mathrm{V}} * ^ {\beta} g
\]

in the sense of the norm \(N_{p}\) , the limit being taken with respect to the section filter of B.

It suffices to apply Prop. 6 (iii), and §2, No.~7, Cor. 3 of Lemma 4.

Remark 3. --- By Prop. 15, the functions \(g * f_{\mathrm{V}}\) , \(f_{\mathrm{V}} * g\) belong to \(\overline{\mathcal{K}(\mathrm{G})}\)

COROLLARY. --- Let W be a closed linear subspace of \(\mathrm{L}^{1}(\mathrm{G},\beta)\) . For W to be a left (resp. right) ideal of \(\mathrm{L}^{1}(\mathrm{G},\beta)\) , it is necessary and sufficient that W be invariant under the left (resp. right) translations of G.

Suppose that W is a left ideal. Let \(s \in G\) and \(g \in W\) . We have \(\varepsilon_{s} * g = \lim_{\mathbb{V}} f_{\mathbb{V}} * (\varepsilon_{s} * g) = \lim_{\mathbb{V}} (f_{\mathbb{V}} * \varepsilon_{s}) * g\) , and \((f_{\mathbb{V}} * \varepsilon_{s}) * g \in \mathbb{W}\) , therefore \(\varepsilon_{s} * g \in W\) , thus \(\boldsymbol{\gamma}(s)g \in \mathbb{W}\) . Conversely, if W is invariant under the left translations, then \(\mu *^{\beta}g \in W\) for \(\mu \in \mathcal{M}^{1}(\mathbf{G})\) and \(g \in W\) , therefore W is a fortiori a left ideal of \(\mathrm{L}^{1}(\mathrm{G}, \beta)\) . One argues similarly for right ideals.

Example. --- We take \(\mathbf{G} = \mathbf{R}\) . Let us define a function \(\mathbf{F}_n \in \mathcal{K}(\mathbf{R})\) by

\[
\begin{array}{l l} F _ {n} (x) = (1 - x ^ {2}) ^ {n} & \text { if } x \in [ - 1, 1 ] \\ F _ {n} (x) = 0 & \text { if } x \notin [ - 1, 1 ]. \end{array}
\]

Let \(\mathbf{A}_n = \int_{-1}^{+1}\mathbf{F}_n(x)dx\) , and \(\mathbf{G}_n = \mathbf{A}_n^{-1}\mathbf{F}_n\) . It is immediate that the measures \(\mathbf{G}_n(x)dx\) satisfy the conditions of §2, No.~7, Cor. 1 of Lemma 4. Let \(\mu\) be a measure on \(\mathbf{R}\) whose support is contained in \([-1/2, 1/2]\) . Then

\[
\begin{array}{r l} & {(\mu * \mathrm{G} _ {n}) (x) = \int_ {\mathbf {R}} \mathrm{G} _ {n} (x - y) d \mu (y)} \\ & {\qquad = \mathrm{A} _ {n} ^ {- 1} \int_ {- 1 / 2} ^ {1 / 2} \mathrm{F} _ {n} (x - y) d \mu (y).} \end{array}
\]

If \(-1 / 2\leqslant x\leqslant 1 / 2\) , then

\[
(\mu * \mathrm{G} _ {n}) (x) = \mathrm{A} _ {n} ^ {- 1} \int_ {- 1 / 2} ^ {1 / 2} [ 1 - (x - y) ^ {2} ] ^ {n} d \mu (y),
\]

therefore \(\mu * \mathbf{G}_n\) coincides on \([-1/2, 1/2]\) with a polynomial. In particular, if \(f\) is a continuous function with support contained in \([-1/2, 1/2]\) , then \(f * \mathbf{G}_n\) coincides in \([-1/2, 1/2]\) with a polynomial; moreover, by Prop. 5 (iv), and §2, No.~7, Cor. 3 of Lemma 4, \(f * \mathbf{G}_n\) converges uniformly to \(f\) . *If \(f\) is of class \(\mathbf{C}^r\) , the derivatives \(\mathrm{D}^s(f * \mathbf{G}_n)\) tend uniformly to \(\mathrm{D}^s f\) for \(0 \leqslant s \leqslant r\) .*

